\section{The converse when \texorpdfstring{$E[2]$}{E[2]} is irreducible}
\label{sec:sa:saturation}

The local matrices now supply the families needed for interpolation.
We shall arrange that a coefficient symbol is a unit at every nonzero
vertex of a binary cube, with all added primes absent at zero.
The key observation is that a unit evaluation has only boundedly many
singular matrix blocks.  A maximal list of such blocks will therefore
force all the required unit evaluations.  In the odd case, a second
construction supplies simultaneous nonvanishing for an auxiliary even
symbol.

\begin{theorem}[The irreducible converse]
\label{thm:irreducible-converse}
Let $E/\Q$ be an elliptic curve for which $E[2]$ is irreducible as a
$G_{\Q}$-module.  If
\[
 c_2(E)=\corank_{\Ztwo}\Sel_{2^\infty}(E/\Q)\in\{0,1\},
\]
then $\ord_{s=1}L(E,s)=c_2(E)$.
\end{theorem}

The curve is fixed throughout.  Every constant may depend on it and on
the finite choices specified below, all made before the cube dimension
grows.  The result applies pointwise to any quadratic twist by taking
that twist as the initial curve.

\subsection{Finite evaluation tables and ordered profiles}

A block is a finite set of the three-state sites of
Section~\ref{sec:lm:local-matrices}, together with their prescribed
matrix entries. A state chooses one of the three alternatives at each
site. Its evaluated matrix $M_s$ deletes the coordinates of absent sites
and applies the toggles of the chosen active alternatives. The empty state has determinant $1$.
We call a state singular when $\det M_s=0$.

The matrix and the coefficient symbol give different information about
such a state. The raw matrix $B$ records the active-prime Selmer equations;
$M$ is the symmetric matrix obtained after the fixed-support correction.
The symbol still depends on the arithmetic realization, even when those
matrix entries have been prescribed. What links them is
Proposition~\ref{prop:selmer-matrix-bound}: a symbol value $1$ permits
only boundedly many singular blocks. We shall first stabilize the
remaining arithmetic dependence and then use this bound.

Fix a normalized coefficient symbol $A$ on the filter of the parameter
$1$. For the odd source retain its condition at the moving detecting
prime $\ell$. Consider an actual finite array of blocks, numbered in
increasing order, with $M=0$ between different blocks. For chosen
indices $i_1<\cdots<i_j$ and states $s_1,\ldots,s_j$, write
\[
 \mathcal E(i_1,\ldots,i_j;s_1,\ldots,s_j)\in\F_2
\]
for its coefficient-symbol value. Each table entry uses its own
aggregate-root probe, which restores the total support label of $1$
and gives the chosen diagonals at active split sites. Root substitution
makes its value independent of the probe with these data. The calibration
below will realize all entries used by a cube with one common set of
linearly activated root primes. The total label, rather than the aggregate root's own label,
is fixed: deleting a site transfers that site's labels to the root.
The substitution rules therefore preserve this value when an unused
site or an empty block is deleted. In particular, the empty table entry
is $A(1)$.

There are two independent indices in these tables. Inside each block,
site slots specify where a state and its addressing paths can be put.
The indices $i_1<\cdots<i_j$ specify ordered copies of the whole block.
We will use a common countable library of site slots in every copy,
closed under the finite addressing constructions of
Lemma~\ref{lem:lm:address-program}. A state in this library uses only
finitely many sites, with every other site absent. A \emph{word} is an ordered finite
list of states from this library, with empty states omitted. An
\emph{ordered profile} assigns a value to each such word, consistently
with the three-state relations in each block. The next lemma constructs
one by stabilizing finite tables: its value depends on the ordered
states, but not on which increasing list of block indices carries them.
It asserts no invariance under permuting the states.

\begin{lemma}[A finitely realizable ordered profile]
\label{lem:sa:profile}
For the fixed symbol $A$, there is a countable library of site slots in each block
and an ordered profile with the following
properties.
\begin{enumerate}[label=\textup{(\roman*)}]
\item Every prescribed finite part is realized by an array of distinct
primes with the data of Lemma~\ref{lem:lm:arithmetic-realization}.
For the odd symbol, it is realized at arbitrarily advanced stages of
the prescribed precision stabilization.
\item A word evaluation depends on the states and on their order, and
not on the particular distinct block indices used to place them.
The tensor relations hold separately in all of its blocks.
\item The site library contains a transferred minimal-weight unit witness and
the primary split sites needed for every finite collection of requests
in
Lemma~\ref{lem:lm:address-program}.
\item For a constant $C$ independent of the finite array, a unit
evaluation contains at most $C$ singular blocks.
\end{enumerate}
Couplings to finitely many fixed primes, kept outside the address
symbol, may be prescribed in advance.  In the even case, all primes
may also satisfy a bound on $v_2(p^2-1)$ independent of the finite
array.
\end{lemma}

We first derive the consequence of this lemma. Its proof follows the
saturation argument and explains how the ordered values retain finite
arithmetic realizations, including for the moving odd source.

\begin{lemma}[Saturation of a punctured cube]
\label{lem:sa:punctured-cube}
For the profile in Lemma~\ref{lem:sa:profile}, either the empty word
has value~$1$, or there is a fixed finite word $W$ of singular blocks
with value~$1$ such that, for every sufficiently large $b$, an addressed
family has symbol~$1$ at every nonzero point of $\F_2^b$ and has all
sites absent at zero.  Every such finite cube is arithmetically
realizable.
\end{lemma}

\begin{proof}
In the tensor space of a block, every nonsingular state is the empty
state plus a sum of singular states, by
Lemma~\ref{lem:binary-tensor-span}.  Expand the transferred unit
witness this way.  Some word consisting entirely of singular blocks,
possibly the empty word, has value~$1$.

The lengths of all such words are bounded by
Lemma~\ref{lem:sa:profile}\textup{(iv)}.  Choose one of maximal
length $m$.  If $m=0$, the empty word has value~$1$.  Otherwise fix
one particular maximizing word
\[
 W=(W_1,\ldots,W_m).
\]
Each $W_i$ uses finitely many slots.  These slots, and thus their total
number, are fixed before $b$ grows.  A bound on the number of blocks
alone would not give this conclusion.

For each $a\in\F_2^b\setminus\{0\}$, place a copy of $W$ in its
original block order.  Address its $i$th block to $a$ using
Lemma~\ref{lem:lm:address-program}.  At $a$ the state is exactly
$W_i$, with every path site absent.  At every other nonzero address
the evaluated block is nonsingular, and at zero it is empty.  All
these finitely many programs are available in the recursive slot list.

At a nonzero address $x$, expand every block outside the copy
designated for $x$ as empty plus singular states.  Choosing the empty
state in every such block leaves $W$ in its original order, hence
gives~$1$.  Every other term contains its $m$ singular blocks and at
least one additional singular block, so maximality forces its value
to be zero.  This includes alternatives involving path slots:
maximality was taken over the whole slot profile.  Multilinearity
therefore gives~$1$ at $x$.  All evaluations use one finite restriction
of the profile, which is realizable by Lemma~\ref{lem:sa:profile}.
\end{proof}

\begin{proof}[Proof of Lemma~\ref{lem:sa:profile}]
\emph{The witness and available slots.}
Begin with the nonroot positions of a minimal-weight unit witness.
The minimum-transfer rule in Theorem~\ref{thm:symbol-rules} moves them
to fresh primes while retaining their fixed-support Frobenius labels.
Neither prescribed internal couplings nor couplings to stationary
primes change this transfer: the latter primes are outside the
coefficient index.  The total simple parity is even, so pair its
simple positions into sites.  Fix admissible generator templates,
including support squareclasses and self types, for this pattern.

At each finite stage place a fresh copy of the pattern in every block
to be colored. First apply Lemma~\ref{tc:minimum-transfer} with
matched actual root probes. The nonroot support labels are unchanged,
so retaining the aggregate root's support label restores the same total
parameter label. Adjust only its bits at active vertices to give the
prescribed diagonals. By Proposition~\ref{tc:irreducible-rules}\textup{(ii)}, a split
root-bit derivative deletes weight one and vanishes below the minimal
weight, while a simple root-bit derivative is zero. Thus every singleton
witness still has value~$1$, even after all internal and cross-block
entries have been prescribed.  For the odd symbol, make the transfer
at the actual coefficient level and retain the witness's labels at
the detecting prime.  Thus the witness entry of every singleton color
is~$1$ before taking homogeneous restrictions.

We now enlarge the site library inside each block.  For every finite list of earlier
site slots and every finite list of addressing requests on it, allocate
fresh primary split site slots realizing all those requests.  Programs
used together receive distinct path slots.  Set every undesignated
coupling to earlier slots equal to zero in $M$.  This gives a
countable formal list, since there are countably many finite requests.
Any finite set of requested slots, together with its finitely many
predecessors, is realizable by
Lemma~\ref{lem:lm:arithmetic-realization}.
The two simple-pair variants have identical outside bits and the
prescribed internal toggle; the aggregate root supplies the split
diagonal alternatives.  Corresponding types have the same reciprocity
correction $L$, so a forward entry $B=M+L$ gives the required reverse
entry as well.  In every alternative impose $M=0$ between blocks.

\emph{Order-invariant evaluations.}
Fix the first $q$ sites of the common library. There are finitely many
states on these sites, and their matrix entries have been prescribed.
For each $j\le q$, color an increasing $j$-tuple of block indices by
the finite table of values
$\mathcal E(i_1,\ldots,i_j;s_1,\ldots,s_j)$ on these states.  There are finitely
many colors.  Apply finite Ramsey successively for these finitely
many arities.  Starting with sufficiently many blocks gives a
homogeneous list of any prescribed finite length.

Let the lengths, arities, and slot cutoffs tend to infinity.
At each fixed cutoff there are only finitely many possible color
vectors.  Successive subsequences and a diagonal subsequence make
every finite vector constant.  These compatible values define the
ordered profile.  Every finite restriction occurs in an original
finite array.  The tensor identities, deletion rules, and unit
witness survive because each involves only such a restriction.

For the odd symbol, carry out this construction along its fixed
nonprincipal ultrafilter of coefficient precisions.  The witness has
bounded weight, so its fixed-$S$ generator squareclasses and binary
self data have only finitely many possibilities; first stabilize these
data. Finer ray-neighborhood depths are realization data and need not
stabilize. Impose the extra labels at $\ell$ at each actual level.  For every fixed weight
bound, Lemma~\ref{cs:fixed-weight-uniformity} gives an
ultrafilter-large set on which all required lower bounds and
smaller-weight vanishing assertions hold.  Otherwise one could choose
a violating index of that bounded weight at each failing level,
contradicting the definitions of the least depth or the least witness
weight.  More explicitly, let $G_q$ be an ultrafilter-large set on
which the witness, cutoff lower bounds, and minimal-weight vanishing
assertions needed for the $q$th finite coloring hold.  Choose
precisions $M_q\to\infty$ in
$\bigcap_{j\le q}G_j$, and apply finite Ramsey there.  Stabilize the
finitely many colors and take the diagonal construction above.  The
extracted subsequence need not itself be ultrafilter-large: every
finite restriction occurs at arbitrarily high precisions, which is
the required conclusion.  We assert no single level valid at all
weight bounds, and choose no infinite array of primes.

\emph{The bound on singular blocks.}
Proposition~\ref{prop:selmer-matrix-bound} and the uniform unit-symbol
bound in Proposition~\ref{prop:unit-symbol-detection} bound the sum
of the block nullities.  Each singular block contributes at least one.
The type alphabet and its correction cost depend on the fixed support
$S$.  The moving place $\ell$ and any bounded stationary environment
contribute only the bounded additional local cost allowed in
Proposition~\ref{prop:selmer-matrix-bound}; they do not enlarge that
alphabet as the array grows.  This proves~\textup{(iv)} in every
finite realization, and therefore in the profile.

For the final even-symbol assertion, choose fixed templates at
sufficient $2$-adic precision for the witness and its simple variants,
and a fixed sufficiently primary template for the path sites.
Refine each rational residue class to a finite precision determining
a finite value of $v_2(p^2-1)$.  Only finitely many templates occur.
Cross-bit conditions at new odd primes do not change these residues.
The same argument will apply to the root primes used below.
\end{proof}

\subsection{Local calibration and the even case}

The preceding argument determines the symbol values.  We next realize
all its local twist classes simultaneously, while keeping each prime's
activation linear in the cube coordinates.

\begin{lemma}[Calibration by root primes]
\label{lem:sa:root-calibration}
The cubes in Lemma~\ref{lem:sa:punctured-cube} admit realizations by
signed squarefree parameters
\[
 h_x=\prod_p(p^*)^{f_p(x)},\qquad x\in\F_2^b,
 \qquad h_0=1,
\]
where each $f_p:\F_2^b\to\F_2$ is linear.  The prescribed local
classes and diagonal choices hold at every assignment.  Moreover,
$h_x$ may be required to be a local square at any prescribed finite
set of stationary primes.  Calibration adds only root primes, whose
invariant dimension $t_p$ is zero.  In the even case, all added primes
may satisfy a bound on $v_2(p^2-1)$ independent of $b$.
\end{lemma}

\begin{proof}
The address and path programs already give a linear activation form
for each nonroot prime.  Both primes of a simple site have the same
form, so their total simple parity vanishes coefficient by coefficient.
For each coordinate $x_j$, add a fresh root prime $q_j$, active when
$x_j=1$.  Choose its rational radical label at the support to cancel
the coefficient of $x_j$ in the total nonroot label.  The paired-simple
condition makes the residual discriminant obstruction to a rotation
label zero.  Lemma~\ref{lem:lm:arithmetic-realization} and root
substitution in Theorem~\ref{thm:symbol-rules} therefore allow these
choices.

The same roots supply the diagonal choices.  At a nonroot split prime
$p$, the unit twist from the other nonroot primes is a linear function
of $x$.  A desired constant diagonal bit on the active set of $p$ is
represented by that constant times $f_p$.  The correction needed is
therefore linear, say $\sum_j c_{pj}x_j$.  Give $q_j$ the Legendre
bit $c_{pj}$ at $p$.  These bits are independently prescribable for
the different $p$; simple positions need no diagonal correction.
At each stationary prime, choose the root bits to cancel the existing
unit twist coefficient by coefficient.  To check compatibility, let $L$ be the compositum of the rational
quadratic fields recording these finite tests. The fresh-prime radical
classes are independent of the support classes by their valuations.
Moreover, $F\cap L=\Q$ for $G=C_3$; for $G=S_3$ it is at most
the residual discriminant field. Thus agreement on the already matched
discriminant coordinate is the only additional condition for a rotation
in $\Gal(F/\Q)$. Chebotarev supplies the required roots.

Using signed prime discriminants gives the prescribed quadratic
characters at every place, including infinity and~$2$. At each $x$,
the nonroot states and their internal data are those of its table entry.
The product of the active root primes has the support label and the
bits at active split primes of the aggregate-root probe in the
corresponding table entry. Its bits at simple primes do not change
the symbol. Root substitution in Theorem~\ref{thm:symbol-rules}
therefore identifies $A(|h_x|)$ with that profile value. The same
comparison includes $x=0$, where the root collection is empty and all
added primes are absent.  Finally, root support labels lie in a
fixed finite set, so fixed rational $2$-adic residue refinements
determine $v_2(q_j^2-1)$.  The extra conditions, ramified at fresh
odd primes, are compatible with these refinements as in
Lemma~\ref{lem:sa:profile}.
\end{proof}

\begin{proof}[Proof of Theorem~\ref{thm:irreducible-converse} when
$c_2(E)=0$]
By Lemma~\ref{pr:invariance}, the functional equation has
sign~$+1$.  We use the even coefficient source on the filter of~$1$.
To obtain its lower bound, choose an auxiliary twist of analytic rank
one, with opposite infinity sign and matching local classes at $2N$,
by nonvanishing in these fixed local classes
(Proposition~\ref{prop:rough-nonvanishing}).  Exclude its prime
support and the exceptional quadratic fields as in
Proposition~\ref{prop:even-coefficients}.  The reverse genus
comparison and Theorem~\ref{thm:heegner-genus-bound} give the lower
bound required there.  Fix this partner, the exclusions, and the
normalized even symbol before choosing the profile or cube dimension.

If the empty word is a unit, Proposition~\ref{prop:unit-symbol-detection}
gives $L(E,1)\ne0$.  Otherwise
Lemmas~\ref{lem:sa:punctured-cube} and~\ref{lem:sa:root-calibration}
give arbitrarily large cubes with unit symbol at all $x\ne0$.
At each such point,
\[
 L(E^{(h_x)},1)\ne0,
 \qquad
 v(\mathcal L(h_x))-2w(h_x)-s(h_x)\le C_A,
\]
where $C_A$ depends only on the fixed normalized symbol.  The local
classes are fixed and the added-prime valuations $v_2(p^2-1)$ are
bounded independently of $b$.  These are precisely the remaining
hypotheses of Theorem~\ref{thm:cyclotomic-interpolation}, which yields
$L(E,1)\ne0$.
\end{proof}

\subsection{A bounded companion for the odd cube}

Assume $c_2(E)=1$. Parity gives functional sign $-1$, so the address
symbol $A$ is odd and its filter is the filter of $1$. Choose its fixed
genus partner of analytic rank zero, with opposite infinity sign and
matching classes at $2N$, by
Proposition~\ref{prop:rough-nonvanishing}. The irreducible genus bound
supplies the coefficient lower estimate. Choose also an even symbol
$A_0$ on a fixed filter of negative odd fundamental discriminants $k$
for which $2N$ splits in $\Q(\sqrt{k})$. It has no condition at the
moving prime $\ell$. Its complementary genus partner has analytic
rank one and is chosen by the same nonvanishing proposition.
Fix these partners, all excluded prime supports and exceptional fields,
and both coefficient normalizations.

If $A(1)=1$, unit-symbol detection already gives the desired analytic
rank. Otherwise the following construction supplies every arithmetic
hypothesis needed for varying-field Heegner interpolation.

\begin{proposition}[An odd cube with a bounded even companion]
\label{prop:sa:bounded-companion}
For the fixed symbols $A,A_0$ just chosen, suppose $A(1)=0$.
There is a constant $C$ such that, for every sufficiently large $b$,
one can find an odd negative fundamental discriminant $k$ and a finite
family
\[
 h_x=\prod_p(p^*)^{f_p(x)},\qquad x\in\F_2^b,
 \qquad f_p:\F_2^b\longrightarrow\F_2\ \text{linear},
\]
with the following properties.
\begin{enumerate}[label=\textup{(\roman*)}]
\item The primes are distinct good odd primes outside the fixed
exclusions; $h_0=1$, and every $h_x$ has the local squareclasses of
$1$ at the fixed support and at infinity. Each $h_x$ is coprime to $k$.
The field $K=\Q(\sqrt{k})$ satisfies the standing Heegner conditions.
\item The prime support of $k$ has cardinality at most $C$. At every
nonroot prime $p\mid\prod_xh_x$, the parameter $k$ is a local square.
At every nonroot prime $q\mid k$, every $h_x$ is a local square.
\item The simultaneous symbol values are
\[
 A_0(|k|)=1,\qquad
 A(|h_x|)=A_0(|kh_x|)=1\quad(x\ne0).
\]
Consequently $E^{(k)}$ has analytic rank zero and $s(k)\le C$.
\end{enumerate}
For the odd source, these assertions mean simultaneous actual tests at
arbitrarily high admissible precisions after $b$ and its finite set of
tests have been fixed. The constant $C$ is independent of $b$, that
precision, and all moving primes. The parameter $k$ may vary with the
finite realization.
\end{proposition}

The construction uses the even symbol twice: at $k$ to control the
base Selmer group over $K$, and at $kh_x$ to make the product have a
simple zero at every nonzero vertex. We first condition $A_0$ by a
fixed singular prefix. Only after this conditioning will its later
values become determinants.

\begin{lemma}[A determinant-conditioned prefix]
\label{lem:sa:prefix}
For the fixed even symbol $A_0$, there are a finite list $a$ of
singular blocks and an environment consisting of at most one simple
vertex with the following property.  Fix these vertices and their
split diagonal entries.  Let $z$ be any further block of tensor sites
whose alternatives all have $B=L$ to every block of $a$.  Choose
aggregate root data preserving the prefix diagonals and the fixed
total label.  Then
\begin{equation}
 A_0(a,z)=\det M_z.
 \label{eq:sa:conditioned-determinant}
\end{equation}
The left side includes the fixed environment.  Its couplings need
only be compatible with the tensor alternatives.
\end{lemma}

\begin{proof}
If the filter of $A_0$ has odd simple parity, retain one simple
vertex of a unit witness as the environment; otherwise take the
environment empty.  The remaining simple positions have even parity,
so pair them into sites and put all variable positions of the witness
in one block.  Allow aggregate roots with the prescribed total class
throughout.

Consider all actual finite configurations with this environment and
a list of singular blocks, with $B=L$ between blocks, at which
$A_0=1$.  This collection, allowing an empty list, is nonempty.
Indeed, if the witness block is nonsingular, expand it as empty plus
singular states by Lemma~\ref{lem:binary-tensor-span}; some resulting
term has value~$1$.  Proposition~\ref{prop:selmer-matrix-bound}
bounds the lengths of these lists, since the environment has fixed
dimension cost.  Choose an actual maximizing list $a$ and fix all
its primes, singular matrices, and diagonal choices.  Only this one
finite list is retained; no bound on the sizes of arbitrary maximizing
blocks is asserted.

For a proposed further block, fixing the prefix and environment gives
a multilinear symbol functional $f$.  It vanishes on every singular
pure state $z$, for otherwise appending $z$ contradicts maximality.
Also $f(\varnothing)=1$, because its aggregate root preserves the
total label and the prefix diagonal data of the maximizing evaluation.
Singular pure states span the kernel of the determinant functional
by Lemma~\ref{lem:binary-tensor-span}, and the empty state has
determinant~$1$.  Hence $f=\det$, proving
Equation~\eqref{eq:sa:conditioned-determinant}.  Maximality here is
over actual singular states.  It requires no alternative choice for
any fixed prefix prime against a later prime.
\end{proof}

\subsubsection{The number of surviving copies}

Fix a nonempty maximal singular unit word
$W=(W_1,\ldots,W_m)$ in an odd profile, as in
Lemma~\ref{lem:sa:punctured-cube}. Write
\[
 M=\bigoplus_{i=1}^mM_i,\qquad
 D=\bigoplus_{i=1}^mD_i,\qquad N=M+D,
\]
where $D_i$ is a sum of site toggles making $M_i+D_i$ invertible.
The matrix $N$ is invertible and $D$ is alternating. Their common
coordinate space $U$ has even dimension, fixed with the word.

\begin{lemma}[Address copies with even total activation]
\label{lem:sa:balanced-copies}
For every $b\ge2$, the address construction for $W$ may be chosen so
that its symbol is $1$ at every nonzero $x$ and all sites are absent
at zero. At a nonzero $x$, eliminating its invertible path sites
leaves the designated copy of $M$ and an odd number of copies of $N$.
The different copies have zero mutual $M$-entries.
\end{lemma}

\begin{proof}
For every $a\ne0$ choose a linear form $l_a$ with
\begin{equation}
 l_a(a)=1,\qquad \sum_{a\ne0}l_a=0.
 \label{eq:sa:activation-parity}
\end{equation}
To obtain these forms, begin with arbitrary choices satisfying the
first condition. Allowed changes at $a$ form its annihilator.
The annihilators of two distinct nonzero addresses are distinct
hyperplanes and span the entire dual space. Changes at those two
addresses can therefore cancel the initial sum.

Activate all word sites in the copy designated for $a$ by $l_a$, and
use a basis of the annihilator of $a$ for its path coordinates in
Lemma~\ref{lem:lm:address-program}. The saturation proof is unchanged.
At a nonzero $x$, the designated copy leaves $M$. Each other active
copy leaves $N$, and a copy with $l_a(x)=0$ leaves only invertible
paths. Equation~\eqref{eq:sa:activation-parity} says that the total
number of active copies is even. The designated copy is active, so
an odd number of other copies remain.
\end{proof}

\subsubsection{A bounded interface}

We will couple every surviving copy to the same fixed collection of
split sites. When two $N$-copies are eliminated, their contributions
to that collection cancel in characteristic two. Thus the calculation
reduces to one $M$-copy, one $N$-copy, and the fixed collection. The
following lemma chooses its couplings; its dimension depends on $W$
and not on the number of addresses.

\begin{lemma}[Alternating interface repair]
\label{lem:sa:interface}
Let $U$ be an even-dimensional vector space over $\F_2$, let $N$
be symmetric and invertible on $U$, and let $D$ be alternating.
There are at most $\dim U$ linear maps $b_i:U\to\F_2^2$ such that,
on writing
$J=\left(\begin{smallmatrix}0&1\\1&0\end{smallmatrix}\right)$ and
$H=\sum_i b_i^tJb_i$, the matrix
\begin{equation}
 S(H)=\begin{pmatrix}D&N\\N&N+H\end{pmatrix}
 \label{eq:sa:interface-matrix}
\end{equation}
is invertible.
If $U$ has the rotation-field structure $\F_4^d$ occurring in the
cyclic cubic case, the maps may be chosen $\F_4$-linear from
$U$ to the field line $\F_4=\F_2^2$.
\end{lemma}

\begin{proof}
Start with $H=0$, and write $\pi_2:U\oplus U\to U$ for projection
to the second coordinate.  For every $H$, this projection is injective
on $\ker S(H)$: if $(u,0)$ belongs to the kernel, then $Nu=0$,
so $u=0$.  The update $H\mapsto H+b^tJb$ adds the alternating
form $(b\pi_2)^tJ(b\pi_2)$ to $S=S(H)$.  We show how to choose
an update reducing every nonzero kernel.

If $\dim\ker S\ge2$, choose $b$ of rank~$2$ on
$\pi_2(\ker S)$.  Take a complement of $\ker S$ and modify its
basis vectors by kernel vectors so that $b\pi_2$ vanishes on the
new complement.  The restriction of $S$ to that complement remains
nondegenerate.  The update is zero there and in its cross pairings,
and has rank~$2$ on $\ker S$.  Hence it decreases the nullity by
exactly~$2$.

Suppose instead that $\ker S=\langle k\rangle$ is one-dimensional.
Let $d_S$ denote the diagonal vector.  Since $z^tSz=d_S^tz$ for
every binary symmetric matrix, $d_S$ annihilates the kernel and lies
in the image of $S$.  Choose $s$ with $Ss=d_S$.  We shall use
\begin{equation}
 d_S^ts=\rank S\pmod2.
 \label{eq:sa:diagonal-rank}
\end{equation}
To prove it, restrict to a nondegenerate complement of the radical,
with invertible symmetric matrix $A$ and diagonal vector $d_A$.
The left side is $d_A^tA^{-1}d_A$.  Off-diagonal terms cancel in
pairs, leaving $\sum_i(A^{-1})_{ii}A_{ii}$.  Applying the same
cancellation to $\tr(A^{-1}A)$ identifies this with $\dim A$
modulo~$2$, proving Equation~\eqref{eq:sa:diagonal-rank}.

Since the ambient dimension is even and the kernel has dimension one,
$d_S^ts=1$.  The alternating matrices $D$ and $H$ have zero
diagonal, so $d_S$ is supported on the second coordinates.
Consequently $\pi_2(s)$ and $\pi_2(k)$ are binary-independent:
otherwise the former would be zero or equal to the latter, forcing
$d_S^ts=0$.  Choose $b$ sending them to independent vectors of
$\F_2^2$.

Let $z$ belong to the kernel after this update.  Self-pairing gives
$d_S^tz=0$.  Pairing the updated equation with $k$ gives
\[
 (b\pi_2(k))^tJ(b\pi_2(z))=0.
\]
Pairing with $s$ gives the same equation with $s$ in place of $k$,
since $s^tSz=d_S^tz=0$.  Their independent images force
$b\pi_2(z)=0$.  The updated equation then reduces to $Sz=0$, so
$z$ is a multiple of $k$.  As $b\pi_2(k)\ne0$, this multiple is
zero.  The update has made $S$ invertible.

Both choices of $b$ remain possible under the field-linear restriction.
Two binary-independent vectors $u,v\in\F_4^d$ are either
field-independent, in which case a field-linear map can send them to
any binary basis of $\F_4$, or satisfy $v=\alpha u$ with
$\alpha\in\F_4\setminus\F_2$.  In the latter case any field-linear
map nonzero on $u$ gives binary-independent images.  These maps are
allowed by the rotation-field off-diagonal spans, and the preceding
argument applies to the underlying binary forms.

Finally, injectivity of $\pi_2$ bounds every kernel dimension by
$\dim U$.  Each update strictly decreases it, so at most $\dim U$
maps are required.
\end{proof}

\begin{proof}[Proof of Proposition~\ref{prop:sa:bounded-companion}]
Choose the prefix and possible one-coordinate environment from
Lemma~\ref{lem:sa:prefix} and fix all their primes. Build the odd
address profile disjoint from them, with $B=L$ to every prefix block.
The alternatives at a simple pair have identical compatible outside
bits against the environment. Include these stationary primes among
the local-square conditions in Lemma~\ref{lem:sa:root-calibration}.

Since $A(1)=0$, Lemma~\ref{lem:sa:punctured-cube} supplies a nonempty
maximal singular unit word $W$. Fix it once. It determines $U,M,D,N$
before $b$ is chosen.

Choose the maps of Lemma~\ref{lem:sa:interface} for this fixed pair
$(N,D)$.  Let $J_{\mathrm{int}}$ be the direct sum of one standard
$J$-block for each map, and let $G$ be the matrix obtained by stacking
the maps $Jb_i$.  Then
\[
 G^tJ_{\mathrm{int}}^{-1}G=\sum_i b_i^tJb_i=H.
\]
For a given $b$, use
Lemma~\ref{lem:sa:balanced-copies} to obtain the address programs and
then calibrate their support labels and diagonals by root primes.

Lemma~\ref{lem:lm:arithmetic-realization} realizes these couplings by
a fixed number of primary split interface positions, with the
field-linear restriction when required.

Only now, for each actual finite address array, construct the auxiliary
parameter $k$.  Its support consists of the stationary prefix, its
possible environment, fresh primes for the fixed number of interface
positions, and one fresh root prime.  Prescribe $B=L$ between the
interface and the prefix blocks.  In the matrix $M$ of the further
block, couple the interface by zero to the path sites and by the same
matrix $G$ to the word coordinates of every address copy.  Keep all
other bits of the simple variants identical.

Require $h_x$ to be a local square at every interface prime for every
$x$.  This is compatible with the matrix entries: the rational norm
bits against address primes are independently prescribable, so their
coefficients against the activation forms can be made zero.  Include
the already chosen address root primes in these conditions; they need
only rational norm bits. Treat them as additional residue places in
Lemma~\ref{lem:lm:arithmetic-realization}. The independent matrix and
norm prescriptions concern the nonroot primes, and the further norm
conditions occur at disjoint root primes. Their quadratic radicals,
ramified at those fresh primes, create no relation with the preceding
extensions, which are unramified there. Root calibration has already made $h_x$ a local
square at the prefix and environment primes.

Choose the last root prime of $k$ to satisfy four requirements: restore
the fixed total support class; make $k$ negative and split at $2N$;
make $k$ a local square at every nonroot address prime; and give the
prescribed diagonals at the prefix and interface.  These are
simultaneous finite rational radical conditions.  Their residual
discriminant coordinate is compatible with a rotation, because the
environment accounts for any odd simple parity of the filter and the
remaining variable simple positions are paired.  The arithmetic
realization lemma therefore supplies the prime.  All new primes can
be taken fresh, so $k$ is squarefree, coprime to every $h_x$, and
avoids the excluded fields.  Address root primes need no splitting
condition, since their invariant dimension is zero.

These choices preserve all matrices needed for the two symbols.
At every actual assignment, the address matrices for $A_0$ equal
those for $A$, since the extra twist by $k$ is a local unit square
at every nonroot address prime.  The prefix and interface matrices
also stay fixed as $x$ varies.  In the tensor evaluations used in
Equation~\eqref{eq:sa:conditioned-determinant}, aggregate roots
preserve these same diagonal choices.

At zero, the further block is just the interface and has invertible
matrix $J_{\mathrm{int}}$.  Equation~\eqref{eq:sa:conditioned-determinant}
therefore gives
\begin{equation}
 A_0(|k|)=1.
 \label{eq:sa:auxiliary-unit}
\end{equation}
At a nonzero address, eliminate the invertible paths.  There remains
one copy of $M$ and an odd number of copies of $N$, each coupled
to the interface by $G$.  Eliminate the $N$-copies in pairs:
each pair has invertible matrix $N\oplus N$ and contributes
$GN^{-1}G^t+GN^{-1}G^t=0$ to the interface.  Retain the one
unpaired copy.  Eliminating the interface then reduces invertibility
of the further block to invertibility of
\[
 \begin{pmatrix}M+H&H\\H&N+H\end{pmatrix}.
\]
The substitution $(u,v)=(u',u'+v')$ transforms this matrix by
congruence to
\[
 \begin{pmatrix}M+N&N\\N&N+H\end{pmatrix}
 =S(H),
\]
which is invertible by Lemma~\ref{lem:sa:interface}.  A second
application of Equation~\eqref{eq:sa:conditioned-determinant} yields
\begin{equation}
 A_0(|k h_x|)=1\quad\text{for every }x\ne0.
 \label{eq:sa:paired-units}
\end{equation}
The odd symbol is also $1$ by saturation. The support of $k$ consists
of the fixed prefix and environment, at most $\dim U$ interface
primes, and one root prime. Its cardinality is bounded independently
of $b$. The symbol $A_0(|k|)=1$ gives analytic rank zero and a uniform
bound for $s(k)$ by Proposition~\ref{prop:unit-symbol-detection}.

For each fixed finite profile restriction, first fix its weight bound
and the finite list of symbol and operator comparisons. Choose a
sufficiently high actual odd-source level, and realize the address
array there as above. Keeping all these address primes fixed, place
the required Galois data in a common finite quotient and choose only
the fresh interface primes and final root prime of $k$ to meet the
remaining matrix and norm prescriptions. All the unit-symbol tests
then hold at once. All constants were fixed
with the symbols, prefix and word. This proves the proposition.
\end{proof}

\subsection{Completion of the odd case}

\begin{proof}[Completion of the proof of
Theorem~\ref{thm:irreducible-converse}]
Apply Proposition~\ref{prop:sa:bounded-companion}. It gives
$\omega(k)$ and $s(k)$ bounded independently of $b$, with
$E^{(k)}$ of analytic rank zero. For $K=\Q(\sqrt{k})$,
Lemma~\ref{lem:h:base-selmer-bound} gives
\[
 \corank_{\Ztwo}\Sel_{2^\infty}(E/K)=1
\]
and bounds the length of its finite quotient by the divisible part
uniformly in $k$.  Its hypotheses hold because $E$ is the fixed
corank-one curve, the other twist has rank zero with bounded $s(k)$,
and $\omega(k)$ is bounded.  This step does not assume finiteness
of the Tate--Shafarevich group of $E$.

For $x\ne0$, the two unit symbols give analytic rank one for
$E^{(h_x)}$ and analytic rank zero for $E^{(kh_x)}$.  Their
coefficient comparisons also give uniform bounds
\[
 j(P_1(h_*))\le w(h_x)+C_1,
 \qquad
 v(\mathcal L(kh_x))\le 2w(kh_x)+C_2.
\]
Here $h_*$ is the fixed genus partner of the odd symbol. The genus point
$P_1(h_*)$ is formed over $\Q(\sqrt{h_xh_*})$, so the point varies with
$x$ while $h_*$ remains fixed. Coprimality
gives $w(kh_x)=w(k)+w(h_x)$. Proposition~\ref{prop:heegner-comparison}
first yields $2j(P_{|h_x|}(h_x))\le4w(h_x)+2w(k)+O(1)$.
Since $w(k)\le\omega(k)$ is uniformly bounded, this becomes
\begin{equation}
 2j(P_{|h_x|}(h_x))\le 4w(h_x)+C_3.
 \label{eq:sa:heegner-upper}
\end{equation}
Subtracting the nonnegative finite lengths $s(h_x)+s(kh_x)$ gives
the uniform upper bound in Theorem~\ref{thm:heegner-interpolation}.

We have also imposed every other hypothesis of that theorem:
$K$ is imaginary and split at $2N$; each $h_x$ is coprime to $k$;
$h_0=1$; signed prime discriminants are activated by linear forms;
the local classes at $2N$ are fixed; and each added prime either
splits in $K$ or is a root prime with $t_p=0$.
For each fixed finite cube, take a sufficiently high precision from
its realization set so that all its unit-symbol tests hold at once.
The constants depend only on the already fixed normalized symbols,
prefix, and word, not on $b$.

The varying-field uniformity in
Theorem~\ref{thm:heegner-interpolation} therefore applies to these
cubes of arbitrarily large dimension.  It gives a simple zero at~$1$
for $L(E,s)L(E^{(k)},s)$.  The second factor is nonzero by
Equation~\eqref{eq:sa:auxiliary-unit}, so $L(E,s)$ has a simple
zero.  This proves the remaining case $c_2(E)=1$.
\end{proof}
